\documentclass[essaypaper,paper]{pmk-graduate}
% Вид документа
% bachelor - Бакалавриат ВКР
% master - Магистратура ВКР
% term - Курсовая работа
% essay - Реферат
%
% assignment - Задание на работу
% paper - Текст работы
%
% thesis - Текст автореферата

\usepackage{amsmath}
\usepackage{multirow}
\usepackage{diagbox}
\usepackage{tabularx}
\usepackage{array}
\usepackage{booktabs}
\usepackage{listingsutf8}

\usepackage{amsmath, amssymb}
\usepackage[ruled, vlined, linesnumbered]{algorithm2e}
\usepackage{hyperref}
\SetKwInput{KwInput}{INPUT}
\SetKwInput{KwOutput}{OUTPUT}
\SetKw{Let}{let}
\SetKw{Exit}{exit}
\SetKwInput{KwInitialize}{Initialize}

\usepackage{tabularx}
\usepackage{array}
\usepackage{booktabs}
\usepackage{ragged2e} % для улучшения разрыва строк
\graphicspath{ {./images/} }

% Библиография оформлена непосредственно в этом файле.
\hypersetup{colorlinks=true,linkcolor=black,citecolor=black,urlcolor=blue}
\begin{document}
% Учебный год
\year{2026-2027}

% Дисциплина
\subject{Информационный менеджмент}

% Студент
\student{Мартиросянц Александр Евгеньевич
\brief А.\,Е.\,Мартиросянц
\group 47
\english Martirosyants Alexander Evgenievich
}
% Научный руководитель
\superviser{Дадеркин Дмитрий Ольгердович
\brief Д.\,О.\,Дадеркин
\position доцент кафедры информатики
\degree д.\,ф.-м.\,н. % Учёная степень. При отсутствии убрать строку
\rank профессор % Учёное звание. При отсутствии убрать строку
}

% Образовательная программа. См. перечень ниже
\program{ics}
% Бакалавриат:
% ics - Информатика и компьютерные науки
% swi - Инженерия программного обеспечения
% swiai - Программная инженерия в искусственном интеллекте
% aida - Искусственный интеллект и анализ данных
% mm - Математическое моделирование
% sa - Системный анализ
% icpim - Инженерное компьютерное проектирование и индустриальная математика
% acse - Прикладная информатика в экономике
% acsm - Прикладная информатика в мехатронике
% icmrs - Интеллектуальное управление в мехатронных и робототехнических системах
% Магистратура:
% macsae - Прикладная информатика в аналитической экономике
% mista - Интеллектуальные системы. Теория и приложения
% memsp - Системное программирование
% memsa - Системный анализ
% mitcd - Информационные технологии в управлении и принятии решений


% Заголовок работы
\title{Связь ERwin с BPwin
    \english The connection between ERwin and BPwin}

% Цель работы
\goal{Исследовать историческую и методологическую связь ERwin и BPwin, их совместное применение и развитие инструментов моделирования.}

% Задачи для разработки в работе
\begin{tasks}
\item Рассмотреть исторические предпосылки IDEF, ERwin и BPwin.
\end{tasks}


% Перечень иллюстративного материала
%\begin{images}
%\image{1}    
%\image{2}    
%\end{images}

% Консультанты
%\begin{advisers}
%\adviser{Сидоров Сидор Сидорович, проектирование схемы БД}   
%\end{advisers}

% Команда для генерации задания на ВКР
%\makeassignment
%\end{document}

% Далее идёт текст самой работы


\maketitlepage

\clearpage

\introduction

Информационный менеджмент рассматривает информацию и данные как ресурсы организации, которыми необходимо управлять на протяжении всего жизненного цикла: от выявления потребностей и описания бизнес-процессов до проектирования, эксплуатации и развития информационных систем. Для согласованной разработки таких систем необходимо описывать как деятельность организации, так и информацию, которую эта деятельность создаёт, использует и изменяет.

\par

В 1990-е — начале 2000-х годов важную роль в практике системного анализа играли CASE-средства (Computer-Aided Software Engineering), предназначенные для построения моделей, документирования требований и автоматизации отдельных этапов проектирования. Среди известных продуктов выделялись BPwin, ориентированный на моделирование бизнес-процессов, и ERwin, предназначенный для моделирования данных и проектирования баз данных. В портфеле Computer Associates эти продукты стали компонентами единого набора средств моделирования.

\par

Актуальность темы состоит в том, что связь между процессами и данными остаётся фундаментальной задачей информационного менеджмента. Хотя конкретные версии BPwin и AllFusion относятся к истории программных средств, вопросы, на которые они отвечали, сохраняют значение: как понять работу организации, определить необходимые данные, устранить противоречия между подразделениями и превратить требования бизнеса в согласованную информационную систему.

\par

\section{Цель и задачи исследования}

Цель работы — исследовать историческую и методологическую связь ERwin и BPwin, причины появления этих инструментов, их совместное применение, дальнейшую эволюцию и современные подходы. Для достижения цели рассматриваются предпосылки развития IDEF; назначения BPwin и ERwin; согласование функциональных и информационных моделей; программные комплексы, поддерживавшие эти подходы; причины изменения рыночного положения продуктов; преимущества и ограничения методологий.

\par

\section{Объект, предмет и методы}

Объект исследования — средства и методологии моделирования в информационном менеджменте. Предмет — взаимодействие BPwin и ERwin при описании бизнес-процессов, данных и требований к информационной системе. Использованы историко-аналитический подход, сопоставление методологий, анализ стандартов и технических материалов производителей.

\par

\chapter{ИСТОРИЧЕСКИЕ ПРЕДПОСЫЛКИ ПОЯВЛЕНИЯ IDEF, ERwin И BPwin}

\section{Проблема сложности информационных систем}

По мере распространения вычислительной техники предприятия столкнулись с тем, что автоматизация отдельных операций не обеспечивала согласованной работы всей организации. Подразделения хранили сведения в несовместимых форматах, использовали разные определения одних и тех же понятий и описывали процессы с неодинаковой степенью подробности. При создании крупной информационной системы требовалось согласовать организационные процедуры, потоки информации, правила обработки и структуру баз данных.

\par

Текстовые описания и неформальные схемы плохо подходили для такой задачи: их было трудно проверять, сравнивать и обновлять. Возникла потребность в формальных графических моделях с определёнными элементами, правилами построения и уровнями детализации. Модель должна была быть понятна специалистам предметной области и одновременно достаточно точна для аналитиков и разработчиков.

\par

Следует различать моделирование процессов и моделирование данных. Первое отвечает на вопросы о функциях, исполнителях, входах и результатах; второе — о сущностях, свойствах, связях и правилах хранения информации. Если разрабатывать эти представления независимо, можно получить процесс, для которого не предусмотрены необходимые данные, или базу данных, не отражающую реальные сценарии работы.

\par

\section{Проект ICAM и формирование семейства IDEF}

Исторические корни IDEF связаны с программой Integrated Computer-Aided Manufacturing (ICAM), которую ВВС США развивали во второй половине 1970-х годов. Целью было повышение производительности промышленности за счёт систематического применения компьютерных технологий. Для интеграции автоматизированных производственных подсистем требовалось единое описание функций и информации, пригодное для использования разными предприятиями и проектными группами.

\par

В рамках ICAM сформировалось семейство методов IDEF — первоначально сокращение от ICAM Definition. IDEF0 предназначался для функционального моделирования, IDEF1 — для информационного моделирования, IDEF2 — для описания динамических характеристик системы. Подходы создавались как взаимодополняющие средства архитектурного описания, а не как конкурирующие способы построения диаграмм.

\par

IDEF0 развивался на основе SADT (Structured Analysis and Design Technique), предложенной Дугласом Россом. В начале 1990-х годов NIST подготовил федеральные стандарты FIPS 183 для IDEF0 и FIPS 184 для IDEF1X. Эти документы закрепили формальные правила и способствовали широкому применению методов в государственном секторе и промышленности [1; 2].

\par

\section{Развитие ERwin и BPwin}

IDEF1X появился как расширение IDEF1 и объединил опыт семантического моделирования данных с принципами реляционной модели. Его назначение — создавать согласованное представление сущностей, атрибутов, ключей и отношений, независимое от особенностей физического хранения. Это позволяло использовать модель как общую основу для проектирования разных приложений и баз данных.

\par

ERwin был разработан компанией Logic Works как средство моделирования данных. Он позволял строить диаграммы сущность—связь, определять атрибуты и ключи, создавать логические и физические модели и генерировать структуры баз данных для поддерживаемых СУБД. В отличие от обычного графического редактора, ERwin сохранял модель как структурированный проект, что позволяло проверять изменения и преобразовывать модель в технические артефакты.

\par

BPwin развивался как CASE-средство функционального моделирования и анализа бизнес-процессов, в первую очередь с поддержкой IDEF0. Он помогал аналитикам фиксировать текущую деятельность предприятия, декомпозировать функции, связывать диаграммы и документировать требования к будущей системе. В конце 1990-х годов продукты перешли в портфель Computer Associates (CA), которая объединила их в линейку ERwin Modeling Suite, а затем AllFusion Modeling Suite. BPwin получил название AllFusion Process Modeler, а ERwin развивался как AllFusion ERwin Data Modeler [3; 4].

\par

Появление пары ERwin—BPwin было ответом на два взаимосвязанных требования: системно описывать работу предприятия и на основе такого описания проектировать информационные структуры. Их связь была методологической — через согласование моделей — и продуктовой — через общий комплекс и механизмы управления моделями.

\par

\chapter{НАЗНАЧЕНИЕ И ПРИНЦИПЫ МОДЕЛИРОВАНИЯ}

\section{BPwin и функциональное моделирование IDEF0}

Основное назначение BPwin — построение функциональных моделей предприятия. В IDEF0 функция изображается прямоугольником, а стрелки показывают взаимодействие функции с окружающей средой. Входы входят слева, выходы выходят справа, управляющие воздействия располагаются сверху, механизмы выполнения — снизу. Это соглашение отделяет преобразуемые объекты от правил, ограничений и ресурсов, обеспечивающих работу.

\par

Модель строится сверху вниз. На контекстной диаграмме A-0 система представлена одной обобщённой функцией и её основными связями с внешним миром. Далее выполняется декомпозиция: функция разбивается на подфункции, раскрывающие детали деятельности. Декомпозиция продолжается до уровня, достаточного для анализа и принятия решений. Между родительской и дочерними диаграммами должны сохраняться согласованные границы и потоки.

\par

IDEF0 удобен для анализа состава работ, информационных входов и результатов. Он позволяет обнаружить дублирование функций, отсутствующие управляющие воздействия, неясные выходы и потенциальные узкие места. Однако IDEF0 не является точной имитационной моделью времени: он показывает функциональную структуру, но не всегда подходит для описания событийной последовательности, временных задержек и параллельного выполнения.

\par

\section{ERwin и моделирование данных IDEF1X}

ERwin предназначен для описания информационных объектов и проектирования схем баз данных. На концептуальном и логическом уровнях аналитик определяет сущности предметной области, атрибуты, идентификаторы, отношения и ограничения. На физическом уровне модель преобразуется в структуру, соответствующую конкретной СУБД: таблицы, столбцы, типы данных, первичные и внешние ключи, индексы и другие элементы.

\par

IDEF1X использует формальные правила для отображения сущностей и отношений. Идентифицирующая связь означает, что ключ родительской сущности участвует в ключе дочерней; неидентифицирующая связь показывает зависимость без включения родительского ключа в первичный ключ дочерней сущности. Модель помогает находить дублирование, проверять ключевые зависимости и создавать согласованные определения данных.

\par

\section{Различия и взаимодополнение моделей}

\small

\setlength{\tabcolsep}{4pt}

\renewcommand{\arraystretch}{1.15}

\begin{longtable}{|p{0.2\textwidth}|p{0.34\textwidth}|p{0.38\textwidth}|}

\hline

\textbf{Критерий} & \textbf{BPwin / IDEF0} & \textbf{ERwin / IDEF1X} \\ \hline

Главный вопрос & Какие функции выполняются? & Какие данные существуют? \\ \hline

Элементы & Функции, входы, выходы, управление, механизмы & Сущности, атрибуты, ключи, отношения \\ \hline

Результат & Функциональная модель & Логическая и физическая модели данных \\ \hline

Пользователи & Бизнес-аналитики, руководители процессов & Архитекторы данных, проектировщики БД \\ \hline

Ограничение & Не описывает подробно схему БД & Не объясняет весь ход бизнес-процесса \\ \hline

\end{longtable}

\normalsize

Модели дополняют друг друга, поскольку любая функция обычно потребляет или создаёт информацию. Например, функция «Оформить заказ» получает сведения о покупателе и товарах, проверяет наличие и формирует заказ. BPwin показывает место функции в работе предприятия и её входы/выходы; ERwin определяет, как представить покупателя, товар, заказ и позиции заказа в базе данных.

\par

Совместное применение даёт более полное описание предметной области. Функциональная модель обеспечивает контекст использования информации, а модель данных проверяет, существует ли структурированный способ эту информацию хранить и связывать. Вместе они позволяют выявлять несогласованность требований до начала программирования.

\par

\chapter{СВЯЗЬ ERwin И BPwin В ЖИЗНЕННОМ ЦИКЛЕ ИС}

\section{Согласование процессов и данных}

Практическая связь ERwin и BPwin заключается не в том, что одна программа заменяет другую, а в согласовании двух представлений одной предметной области. Функциональная модель отвечает за организационный смысл работы, а модель данных — за формализованное представление информации, с которой работают функции. Согласование помогает проверить полноту требований и сделать переход от обследования предприятия к проектированию системы более управляемым.

\par

Типовой процесс начинается с обследования организации и построения модели «как есть» (AS-IS). Аналитик фиксирует функции, входы, выходы, правила и механизмы. Затем выявляются проблемы: ручной повторный ввод, дублирование согласований, отсутствие владельца данных или неясные источники информации. После этого формируется целевая модель «как должно быть» (TO-BE), отражающая предлагаемое улучшение.

\par

На основе процессов определяются информационные потребности: какие сущности должны существовать, какие атрибуты требуются для выполнения функций, какие данные поступают извне и какие результаты необходимо сохранять. Эти требования передаются в ERwin для построения логической модели. Затем проводится обратная проверка: может ли каждая функция получить необходимые данные и не создаются ли данные без понятного потребителя.

\par

\section{Пример комплексного применения}

Рассмотрим интернет-магазин. В BPwin процесс «Обработать заказ» можно разделить на получение заказа, проверку покупателя, проверку наличия товара, расчёт стоимости, подтверждение оплаты и передачу заказа на доставку. Входами будут данные заказа и сведения о покупателе; управляющими воздействиями — правила продажи, тарифы и условия оплаты; механизмами — сотрудники, информационная система и платёжный сервис; выходами — подтверждённый заказ и документы для доставки.

\par

В ERwin эта деятельность отражается сущностями CUSTOMER, ORDER, ORDER\_ITEM, PRODUCT, PAYMENT и SHIPMENT. CUSTOMER хранит идентификатор и контактные данные покупателя. ORDER содержит дату, статус и ссылку на покупателя. ORDER\_ITEM связывает заказ с товарами и фиксирует количество и цену на момент покупки. PAYMENT и SHIPMENT позволяют отдельно учитывать платёжные и логистические события.

\par

\section{Проверка согласованности и CASE-средства}

Сопоставление моделей выявляет важные вопросы. Если процесс допускает повторную оплату или возврат, следует определить, может ли модель данных хранить несколько платёжных операций для одного заказа. Если цена фиксируется в момент покупки, нельзя полагаться только на текущую цену в PRODUCT: необходимо хранить цену в ORDER\_ITEM. Если доставка может быть разделена на несколько отправлений, модель должна поддерживать соответствующую кратность связей.

\par

CASE-средства поддерживают разработку систем через формализованные модели, репозитории, автоматическую генерацию артефактов и проверки. В составе AllFusion Modeling Suite Process Modeler, Model Manager и ERwin могли использоваться совместно. Model Manager обеспечивал библиотеку моделей, управление доступом и версиями, а также синхронизацию моделей данных. CA позиционировала эти инструменты как средства согласования бизнес-требований и ИТ-решений [3; 4].

\par

Общий репозиторий особенно важен, когда над проектом работают несколько групп. Изменение бизнес-правила может повлиять на процесс, структуру данных, отчёты и интерфейсы. Централизованное управление помогает определить актуальную версию, автора изменения и артефакты, требующие пересмотра. Однако наличие общего пакета не гарантирует автоматической семантической синхронизации: необходимы правила именования, договорённости команды и регулярные проверки.

\par

\chapter{ПРОГРАММНЫЕ КОМПЛЕКСЫ И АЛЬТЕРНАТИВЫ}

\section{AllFusion Modeling Suite}

В начале 2000-х годов CA продвигала AllFusion Modeling Suite как интегрированный набор CASE-инструментов. В него входили AllFusion Process Modeler (ранее BPwin), AllFusion ERwin Data Modeler, средства управления моделями и дополнительные компоненты для проверки и публикации. Такой подход соответствовал идее общего цикла: описать бизнес-функции, спроектировать данные, проверить модель и организовать коллективную работу.

\par

Переименование продуктов отражало стремление представить отдельные инструменты как компоненты единой архитектуры моделирования. Пользователи продолжали использовать привычные названия BPwin и ERwin, поэтому в технической и учебной литературе встречаются оба обозначения. В 2006 году CA объявляла выпуск AllFusion Modeling Suite r7.1; позднее линейка продвигалась под брендом CA ERwin Modeling [4; 5].

\par

Пакет не был единственной возможной реализацией методологий. IDEF0 и IDEF1X — это методы и нотации, которые могут поддерживаться разными программами. При выборе инструмента важно оценивать наличие проверок правил, обмен данными, совместную работу, документацию, интеграцию и стоимость владения, а не только возможность нарисовать диаграмму.

\par

\section{Современные средства моделирования процессов}

В современных проектах широко применяется BPMN (Business Process Model and Notation). В отличие от IDEF0, который делает акцент на функциональной декомпозиции, BPMN позволяет отображать последовательности событий, ветвления, сообщения, участников, параллельные действия и взаимодействие процессов. Спецификация BPMN разработана Object Management Group; версия 2.0 принята в 2010 году, а BPMN оформлена как международный стандарт ISO/IEC 19510 [6; 7].

\par

BPMN полезна, когда требуется описать рабочий поток достаточно точно для согласования с владельцами процесса и потенциальной автоматизации. IDEF0 остаётся удобным для обзора функций организации. Следовательно, BPMN не следует считать прямой заменой IDEF0 во всех задачах: выбор зависит от цели модели, требуемой строгости и аудитории.

\par

Программные средства включают специализированные BPM-платформы и универсальные редакторы диаграмм. Microsoft Visio предлагает шаблоны общих блок-схем и методологических диаграмм, включая IDEF0 и BPMN [8]. При выборе необходимо проверить поддержку конкретной версии стандарта и возможности экспорта.

\par

\section{Современные средства моделирования данных}

ERwin продолжил существование как самостоятельное направление моделирования данных. В современных продуктах erwin Data Modeler акцент делается на построении и сопровождении моделей, работе с платформами хранения, документировании и управлении метаданными. Смена владельцев и брендов не означает исчезновения самой задачи моделирования данных: она сохраняет значение при проектировании реляционных баз, хранилищ и корпоративной архитектуры [9].

\par

Кроме ERwin, применяются инструменты, встроенные в среды разработки и СУБД, универсальные средства архитектурного моделирования и редакторы ER-диаграмм. В отдельных проектах схемы описываются как код — например, через миграции базы данных и средства контроля версий. Это удобно для командной разработки, однако не всегда заменяет концептуальную и логическую модель предметной области.

\par

\small

\setlength{\tabcolsep}{4pt}

\renewcommand{\arraystretch}{1.15}

\begin{longtable}{|p{0.2\textwidth}|p{0.34\textwidth}|p{0.38\textwidth}|}

\hline

\textbf{Задача} & \textbf{Тип инструмента} & \textbf{Что оценивать} \\ \hline

Обзор функций & IDEF0-инструмент & Декомпозиция и правила диаграмм \\ \hline

Детальный поток работ & BPMN-редактор / BPM-платформа & События, шлюзы, экспорт \\ \hline

Схема базы данных & ERwin или другой data-modeler & СУБД, генерация SQL, сравнение схем \\ \hline

Управление архитектурой & Репозиторий моделей & Версии, права, трассируемость \\ \hline

Командное изменение схемы & Миграции и контроль версий & Повторяемость развёртывания \\ \hline

\end{longtable}

\normalsize

Современная экосистема часто состоит из нескольких специализированных средств. Это повышает гибкость, но требует отдельно решать вопрос согласованности артефактов и обмена данными между программами.

\par

\chapter{ПРИЧИНЫ ТРАНСФОРМАЦИИ РЫНКА}

Формулировка «BPwin ушёл с рынка» требует уточнения. BPwin не исчез одномоментно: он был переименован в AllFusion Process Modeler, затем существовал в линейке CA ERwin Modeling. CA публиковала обновления Process Modeler и объявления о прекращении поддержки отдельных старых версий. В 2003 году компания сообщала о версии AllFusion Process Modeler 4.1 и завершении поддержки релиза 4.0 в 2004 году [3]. Исторические объявления о выпусках и жизненном цикле сохранились в каталоге поддержки Broadcom [5].

\par

\section{Изменение задач бизнеса}

Одной из причин изменения положения классических CASE-средств стало развитие корпоративных платформ, в которых моделирование процессов связывалось с управлением исполнением, мониторингом показателей и автоматизацией рабочих потоков. От инструмента, который прежде создавал диаграммы и документацию, организации стали ожидать взаимодействия с workflow-системами, сервисами, аналитикой и облачной инфраструктурой.

\par

Второй фактор — распространение BPMN. Общая нотация облегчила обмен процессными моделями между специалистами и программными платформами. BPMN поддерживает формальное представление событий, ветвлений и взаимодействий, что делает её привлекательной для задач, где требуется перейти от схемы к автоматизированному исполнению [6; 7].

\par

Третий фактор — изменение архитектуры разработки. Схемы баз данных могут управляться через миграции, репозитории кода, CI/CD-процессы и инструменты наблюдаемости. Это не устраняет потребность в моделях, но меняет способы их создания и сопровождения. Для некоторых проектов важнее интеграция с конвейером разработки, чем отдельная настольная программа.

\par

\section{Что пришло на смену прежнему комплексу}

Вместо одного интегрированного настольного пакета организации нередко используют набор средств: BPMN-редактор или BPM-платформу для процессов, средство моделирования данных, репозиторий для совместной работы, систему управления требованиями и инструменты контроля версий. Преимущество — выбор инструмента под задачу и интеграция в технологическую среду. Недостаток — необходимость самостоятельно обеспечивать согласованность артефактов.

\par

Методология и программный продукт имеют разные жизненные циклы. Программа может потерять актуальность из-за прекращения поддержки, изменения бизнес-модели производителя, ограниченной совместимости или конкуренции. Нотация при этом может продолжать использоваться, если решает определённую задачу и поддерживается другими инструментами. IDEF0 остаётся полезным для функционального описания, а ER-моделирование — для анализа структуры данных.

\par

Поэтому некорректно утверждать, что BPMN полностью заменила IDEF0 или современные редакторы полностью заменили ERwin. Правильнее говорить о перераспределении ролей: BPMN стала распространённым стандартом детального процессного моделирования; IDEF0 сохраняет нишу функционального анализа; ERwin развивается в направлении моделирования данных; другие инструменты закрывают задачи управления архитектурой и DevOps.

\par

\chapter{ДИНАМИКА РАЗВИТИЯ И ПЕРСПЕКТИВЫ}

\small

\setlength{\tabcolsep}{4pt}

\renewcommand{\arraystretch}{1.15}

\begin{longtable}{|p{0.2\textwidth}|p{0.34\textwidth}|p{0.38\textwidth}|}

\hline

\textbf{Период} & \textbf{Событие / тенденция} & \textbf{Значение} \\ \hline

1970-е — начало 1980-х & ICAM и формирование IDEF & Общий язык описания функций и информации \\ \hline

1980-е & Развитие IDEF1X & Формализация моделирования данных \\ \hline

1993 & NIST FIPS 183 и FIPS 184 & Стандартизация IDEF0 и IDEF1X \\ \hline

1990-е & Распространение CASE-средств & Автоматизация диаграмм и проектирования \\ \hline

Конец 1990-х — начало 2000-х & Переход продуктов к CA и AllFusion & Интеграция моделирования и управления моделями \\ \hline

2000-е — 2010-е & Распространение BPMN & Стандартизация детальных процессов \\ \hline

2010-е — настоящее время & Облака, API, DevOps и data governance & Распределённые инструменты и метаданные \\ \hline

\end{longtable}

\normalsize

\section{Современные направления}

Первое направление — переход от статической документации к моделям, включённым в жизненный цикл изменений. Схемы должны обновляться вместе с требованиями, программным кодом и структурой базы данных. Важны версионирование, автоматические проверки и трассируемость: возможность установить, какое требование породило функцию и какие данные необходимы для её реализации.

\par

Второе направление — интеграция бизнес-архитектуры и управления данными. Организациям недостаточно знать, какие таблицы существуют; необходимо понимать происхождение данных, владельцев, качество, правила использования и связь с бизнес-целями. Модели становятся источником корпоративных метаданных, а не только иллюстрациями в отчёте.

\par

Третье направление — совместное применение нескольких нотаций. IDEF0 подходит для иерархического описания функций; BPMN — для детального потока работ; ER-модели — для структуры данных; UML — для проектирования программных компонентов. Выбор одной универсальной нотации редко оптимален для всех уровней абстракции.

\par

\section{Перспективы IDEF0 и IDEF1X}

IDEF0 сохраняет полезность для анализа деятельности, регламентирования, функциональной декомпозиции и обучения системному мышлению. Его сильная сторона — представление сложной организации в виде набора функций и связей без излишних технических подробностей. Для описания динамики и событийных сценариев его обычно дополняют другими средствами.

\par

IDEF1X сохраняет ценность там, где нужна строгая модель сущностей и отношений, особенно в реляционных системах и проектах с высокими требованиями к документации. Современные технологии не отменяют необходимости определять ключи, ограничения целостности и семантику атрибутов. Развитие идёт не к отказу от моделирования, а к связанным системам моделей и метаданных.

\par

\chapter{ПРЕИМУЩЕСТВА И НЕДОСТАТКИ ПОДХОДА ERwin—BPwin}

\section{Преимущества}

\noindent 1. Комплементарность представлений: BPwin раскрывает функции и потоки, ERwin — информационные объекты и отношения.\par

\noindent 2. Раннее выявление несоответствий: сопоставление функций с необходимыми данными помогает обнаружить пропущенные атрибуты и неопределённые источники.\par

\noindent 3. Формализация требований: графические модели создают общую основу для общения заказчиков, аналитиков, архитекторов и разработчиков.\par

\noindent 4. Повторное использование и документирование: модели можно применять при разработке нескольких приложений, обучении и аудите.\par

\noindent 5. Переход к реализации: ERwin поддерживал генерацию физических схем, а BPwin давал структурированное описание функций.\par

\section{Недостатки и ограничения}

\noindent 1. Стоимость внедрения и обучения: формальные методологии требуют подготовки специалистов и единых правил.\par

\noindent 2. Риск устаревания моделей: если процессы и база меняются, а модели нет, документация становится недостоверной.\par

\noindent 3. Ограниченность каждой нотации: IDEF0 не является языком исполняемых процессов, а ER-модель не описывает всю бизнес-логику.\par

\noindent 4. Интеграционные трудности: разные инструменты могут использовать различные форматы и понятия, а экспорт не всегда сохраняет смысл связей.\par

\noindent 5. Зависимость от жизненного цикла продукта: смена владельца, прекращение поддержки и несовместимость форматов могут затруднить сопровождение.\par

\section{Сводная оценка}

\small

\setlength{\tabcolsep}{4pt}

\renewcommand{\arraystretch}{1.15}

\begin{longtable}{|p{0.2\textwidth}|p{0.34\textwidth}|p{0.38\textwidth}|}

\hline

\textbf{Критерий} & \textbf{Сильная сторона} & \textbf{Риск} \\ \hline

Методология & Разделение функций и данных & Нужна дисциплина согласования \\ \hline

Автоматизация & Генерация схем и документации & Не гарантирует правильность требований \\ \hline

Командная работа & Общие модели и версии & Конфликты и интеграционные сложности \\ \hline

Сопровождение & Сохранение знаний о системе & Нужно обновлять артефакты \\ \hline

Современная применимость & Принципы остаются актуальными & Старые версии могут не поддерживаться \\ \hline

\end{longtable}

\normalsize

\conclusion

Исследование показывает, что ERwin и BPwin возникли как практический ответ на потребность описывать предприятие с нескольких взаимосвязанных сторон. Историческое развитие IDEF в рамках программы ICAM было направлено на создание формализованного языка для анализа функций, информации и динамики промышленных систем. Позднее IDEF0 и IDEF1X получили стандартизацию, что способствовало распространению соответствующих подходов.

\par

BPwin обеспечивал функциональное моделирование, позволяя представить деятельность предприятия в виде иерархии функций, входов, выходов, управляющих воздействий и механизмов. ERwin поддерживал проектирование данных: определение сущностей, атрибутов, ключей и отношений, а также переход от логической модели к физической схеме базы данных. Их связь проявлялась в том, что функции предприятия нуждаются в информации, а структура данных должна обеспечивать выполнение реальных бизнес-сценариев.

\par

В составе Computer Associates оба направления вошли в AllFusion Modeling Suite, включавший Process Modeler, ERwin Data Modeler и инструменты управления моделями. Конкретные продукты переименовывались и проходили этапы изменения поддержки. Поэтому говорить о полном исчезновении методологий некорректно: изменился рынок инструментов, но задачи функционального и информационного моделирования сохранились.

\par

Распространение BPMN, развитие BPM-платформ, облачных сервисов, репозиториев и DevOps привели к тому, что вместо одного настольного комплекса всё чаще применяется набор специализированных средств. Современный подход требует поддерживать согласованность процессов, требований, данных, программной реализации и правил управления изменениями. Историческая значимость ERwin и BPwin состоит в демонстрации принципа комплексного моделирования: система должна отражать и деятельность организации, и данные, необходимые для её выполнения.

\par

\chapter*{Выводы}

\addcontentsline{toc}{chapter}{Выводы}

\noindent 1. ERwin и BPwin решали разные, но взаимосвязанные задачи: моделирование данных и функциональное моделирование предприятия.\par

\noindent 2. Исторические предпосылки их методологической связи сформировались в рамках программы ICAM и семейства IDEF.\par

\noindent 3. IDEF0 описывает функции, входы, выходы, управление и механизмы; IDEF1X формализует сущности, атрибуты, ключи и отношения.\par

\noindent 4. Совместное применение позволяет проверять, обеспечивают ли проектируемые данные реальные бизнес-процессы.\par

\noindent 5. AllFusion Modeling Suite объединил Process Modeler, ERwin Data Modeler и средства управления моделями.\par

\noindent 6. BPwin не исчез одномоментно: продукт переименовывался и развивался как AllFusion Process Modeler; менялись версии и жизненный цикл поддержки.\par

\noindent 7. BPMN стала важным стандартом детального моделирования процессов, но не является универсальной заменой IDEF0.\par

\noindent 8. ER-моделирование остаётся значимым для проектирования и сопровождения баз данных, хотя инструменты разработки изменились.\par

\noindent 9. Преимущества подхода — согласованность, формализация требований, документирование и раннее выявление ошибок.\par

\noindent 10. Ограничения — стоимость обучения, риск устаревания моделей, ограниченность нотаций и сложности интеграции.\par

\noindent 11. Современным проектам нужна трассируемость между процессами, требованиями, моделями данных и реализацией.\par

\noindent 12. Методологии могут сохранять актуальность значительно дольше, чем отдельные программы, которые их поддерживают.\par

\chapter*{СПИСОК ИСПОЛЬЗОВАННЫХ ИСТОЧНИКОВ}

\addcontentsline{toc}{chapter}{СПИСОК ИСПОЛЬЗОВАННЫХ ИСТОЧНИКОВ}

1. National Institute of Standards and Technology. Integration Definition for Function Modeling (IDEF0): Federal Information Processing Standards Publication 183. — Gaithersburg : NIST, 1993. — URL: \url{https://nvlpubs.nist.gov/nistpubs/Legacy/FIPS/fipspub183.pdf} (дата обращения: 25.10.2026).

\par

2. National Institute of Standards and Technology. Integration Definition for Information Modeling (IDEF1X): Federal Information Processing Standards Publication 184. — Gaithersburg : NIST, 1993. — URL: \url{https://nvlpubs.nist.gov/nistpubs/Legacy/FIPS/fipspub184.pdf} (дата обращения: 25.10.2026).

\par

\noindent 3. Computer Associates International, Inc. AllFusion Process Modeler 4.0 Product Support Announcement. — 2003. — URL: \url{https://ftpdocs.broadcom.com/phpdocs/0/260/BPwin4.0\_PSA.pdf} (дата обращения: 25.10.2026).\par

\noindent 4. CA представляет новые возможности AllFusion ERwin Data Modeler (ERwin), AllFusion Process Modeler (BPwin) и других средств моделирования : пресс-релиз от 08.12.2003. — URL: \url{https://www.interface.ru/ca/news/pr031208715.htm} (дата обращения: 25.10.2026).\par

5. Broadcom. CA ERwin Modeling Suite Product Release and Support Announcements. — URL: \url{https://support.broadcom.com/web/ecx/support-content-notification/-/external/content/release-announcements/CA-ERwin-Modeling-Suite-Product-Release-and-Support-Announcements/4551} (дата обращения: 25.10.2026).

\par

\noindent 6. Object Management Group. Business Process Model and Notation (BPMN), Version 2.0. — 2010. — URL: \url{https://www.omg.org/spec/BPMN/2.0/PDF} (дата обращения: 25.10.2026).\par

\noindent 7. International Organization for Standardization. ISO/IEC 19510:2013. Information technology — Object Management Group Business Process Model and Notation. — Geneva : ISO, 2013. — URL: \url{https://www.iso.org/standard/62652.html} (дата обращения: 25.10.2026).\par

\noindent 8. Microsoft. Process diagrams in Visio. — URL: \url{https://support.microsoft.com/en-us/visio/process-diagrams-in-visio} (дата обращения: 25.10.2026).\par

\noindent 9. Маклаков С. В. Создание информационных систем с AllFusion Modeling Suite. — М. : Диалог-МИФИ, 2003. — URL: \url{https://djvu.online/file/K1W6pNq9dnRxi} (дата обращения: 25.10.2026).\par

\noindent 10. erwin, Inc. About Us: The Evolution — From Data Modeling to Data Governance. — URL: \url{https://bookshelf.quest.com/about-us/} (дата обращения: 25.10.2026).\par

\noindent 11. National Institute of Standards and Technology. Integration Definition for Information Modeling (IDEF1X). — 1993. — URL: \url{https://www.govinfo.gov/app/details/GOVPUB-C13-986bf8b12a4fed44eb78fca0bb55d668} (дата обращения: 25.10.2026).\par

\end{document}
